% Image credits for the photographs used in this volume (Book 1), Dutch
% edition. Machine-readable license records live in images/book1/CREDITS.md;
% keep this file in sync with frontmatter/image-credits.tex.
\chapter*{Beeldverantwoording}

De tekeningen in dit boek zijn origineel. De foto's worden gebruikt onder
hun respectieve vrije licenties, met dank aan hun makers:

\begin{itemize}
  \item \emph{John Dalton} (hoofdstuk over atomen en moleculen): gravure
        door Charles Turner naar Joseph Allen, 1834, Library of Congress,
        via Wikimedia Commons, publiek domein.
  \item \emph{Daltons tabel van de elementen, 1808} (hoofdstuk over atomen
        en moleculen): plaat uit John Daltons \emph{A New System of
        Chemical Philosophy}, scan van het Internet Archive, via Wikimedia
        Commons, publiek domein.
  \item \emph{Bauxiet} (hoofdstuk over grondstoffen): foto door Werner
        Schellmann, via Wikimedia Commons, CC BY-SA 2.5.
  \item \emph{Graniet} (hoofdstuk over zuivere stoffen): foto door Eurico
        Zimbres, via Wikimedia Commons, CC BY-SA 2.0 br.
  \item \emph{Watervrij kopersulfaat} (hoofdstuk over stoffen herkennen):
        foto door W.~Oelen, via Wikimedia Commons, CC BY-SA 3.0.
  \item \emph{Kopersulfaatkristallen} (hoofdstuk over stoffen herkennen):
        foto door Stephanb, via Wikimedia Commons, CC BY-SA 3.0.
  \item \emph{Antoine-Laurent Lavoisier en Marie-Anne Paulze Lavoisier}
        (hoofdstuk over kloppende reactievergelijkingen): schilderij van
        Jacques-Louis David, 1788, Metropolitan Museum of Art, via
        Wikimedia Commons, publiek domein.
  \item \emph{Bakelieten telefoon} (hoofdstuk over kunststoffen): foto
        door Holger Ellgaard, via Wikimedia Commons, CC BY-SA 3.0.
  \item \emph{Ernest Rutherford} (hoofdstuk over het atoom): foto, George
        Grantham Bain Collection, Library of Congress, via Wikimedia
        Commons, publiek domein.
  \item \emph{Haliet} (hoofdstuk over ionen): foto door Robert M.~Lavinsky,
        via Wikimedia Commons, CC BY-SA 3.0.
  \item \emph{Het Vrijheidsbeeld} (hoofdstuk over metalen en corrosie):
        foto door Elcobbola, via Wikimedia Commons, publiek domein.
  \item \emph{Dmitri Mendelejev} en \emph{zijn tabel van 1869} (hoofdstuk
        over het periodiek systeem): via Wikimedia Commons, publiek domein.
  \item \emph{De Krabnevel} (hoofdstuk over de elementen): NASA, ESA,
        J.~Hester en A.~Loll, publiek domein.
  \item \emph{Natrium} (hoofdstuk over elektronenschillen): foto door
        Dnn87, via Wikimedia Commons, CC BY-SA 3.0.
  \item \emph{Chloor in een ampul} (hoofdstuk over elektronenschillen):
        foto door W.~Oelen, via Wikimedia Commons, CC BY-SA 3.0.
  \item \emph{Amedeo Avogadro} (hoofdstuk over de mol): lithografie naar
        een tekening van C.~Sentier, 1856, Edgar Fahs Smith collection, via
        Wikimedia Commons, publiek domein.
  \item \emph{Sneeuwkristallen} (hoofdstuk over polariteit): foto's door
        Wilson Bentley, 1902, via Wikimedia Commons, publiek domein.
  \item \emph{Louis Pasteur} (hoofdstuk over stereochemie): foto door Paul
        Nadar, via Wikimedia Commons, publiek domein.
  \item \emph{Kristallen van ammoniumtartraat} (hoofdstuk over
        stereochemie): foto door Marie-Lan Taÿ Pamart, via Wikimedia
        Commons, CC BY 4.0.
  \item \emph{Svante Arrhenius} (hoofdstuk over zuren en basen):
        fotogravure, 1909, via Wikimedia Commons, publiek domein.
  \item \emph{Een zuil van Volta} (hoofdstuk over elektrochemische
        cellen): foto door GuidoB, via Wikimedia Commons, CC BY-SA 3.0.
  \item \emph{Wallace Carothers in zijn laboratorium} (hoofdstuk over
        polymeren): fotograaf onbekend, via Wikimedia Commons, publiek
        domein.
\end{itemize}

De volledige bronlinks en licentie-URL's van elke foto staan in
\texttt{images/book1/CREDITS.md} in de repository van het boek.
De gevarenpictogrammen zijn die van het Globally Harmonized System of
Classification and Labelling of Chemicals, zoals getekend voor het pakket
\texttt{ghsystem} van \TeX\ Live.

\bigskip
Naast de foto's en de originele tekeningen zijn sommige figuren
illustraties die met kunstmatige intelligentie zijn gemaakt, met de
beeldgeneratie van OpenAI en op basis van prompts die de redacteuren van
het boek hebben geschreven, en die vóór opname op chemische juistheid zijn
nagekeken. De volledige prompt van elk gegenereerd beeld staat in
\texttt{images/book1/ai/PROMPTS.md} in de repository.
\clearpage
